% Kapitel 5: Evaluation
\section{Evaluation}
\label{sec:evaluation}

\subsection{Methode}

Die Evaluation erfolgte durch Nutzertests mit fünf Probanden (N=5) aus der Zielgruppe Studierende. Die Tests folgten einem Think-Aloud-Protokoll mit drei Aufgaben:

\begin{enumerate}
    \item \textbf{Aufgabe 1:} ``Finde heraus, welche VR-Projekte das Lab gerade betreibt.''
    \item \textbf{Aufgabe 2:} ``Kannst du herausfinden, ob du als Student Geräte ausleihen kannst?''
    \item \textbf{Aufgabe 3:} ``Finde eine Publikation zum Thema EEG oder Biosensoren.''
\end{enumerate}

Nach den Aufgaben füllten die Probanden den SUS-Fragebogen aus. Zusätzlich wurde der Chatbot gegen die 18 Amershi-Guidelines evaluiert.

\subsection{SUS-Ergebnisse}

\begin{tabularx}{\textwidth}{lX}
\toprule
\textbf{Proband} & \textbf{SUS-Score} \\
\midrule
P1 & 72,5 \\
P2 & 65,0 \\
P3 & 77,5 \\
P4 & 70,0 \\
P5 & 82,5 \\
\midrule
\textbf{Mittelwert} & \textbf{73,5} \\
\bottomrule
\end{tabularx}

Mit einem Mittelwert von 73,5 liegt Ergonomia über dem Benchmark von 68~\cite{lewis2018} und damit im Bereich ``gut'' (B-Note nach der SUS-Schulnoten-Skala).

\subsection{Amershi-Evaluation}

\begin{tabularx}{\textwidth}{lXX}
\toprule
\textbf{\#} & \textbf{Guideline} & \textbf{Erfüllt} \\
\midrule
G1 & Make clear what the system can do & \checkmark \\
G2 & Make clear how well the system can do what it can do & \checkmark \\
G3 & Time services based on context & \checkmark \\
G4 & Show contextually relevant information & \checkmark \\
G5 & Match relevant social norms & \checkmark \\
G6 & Mitigate social biases & -- \\
G7 & Support efficient invocation & \checkmark \\
G8 & Support efficient dismissal & \checkmark \\
G9 & Support efficient correction & -- \\
G10 & Scope services when in doubt & \checkmark \\
G11 & Make clear why the system did what it did & \checkmark \\
G12 & Remember recent interactions & \checkmark \\
G13 & Learn from user behavior & -- \\
G14 & Update and adapt cautiously & -- \\
G15 & Encourage granular feedback & -- \\
G16 & Convey consequences of user actions & -- \\
G17 & Provide global controls & -- \\
G18 & Notify users about changes & -- \\
\bottomrule
\end{tabularx}

Die Evaluation zeigt, dass Ergonomia die zentralen Guidelines für den initialen Use Case gut erfüllt, während fortgeschrittene Features (adaptives Lernen, Feedback-Mechanismen) für das Projektziel nicht erforderlich waren.

\subsection{Qualitative Befunde}

Die Think-Aloud-Protokolle ergaben folgende wiederkehrende Themen:

\begin{itemize}
    \item Positiv: ``Die Antworten sind präzise und die Quellenangaben helfen, die Ergebnisse zu verifizieren.''
    \item Positiv: ``Der Chatbot ist schnell und die Antworten kommen flüssig.''
    \item Verbesserung: ``Einige Fachbegriffe werden nicht erkannt -- da wäre ein Synonym-Feature hilfreich.''
    \item Verbesserung: ``Vorschlagsfragen zum Einstieg würden den Nutzungseinstieg erleichtern.''
\end{itemize}
